\chapter{绪论}

\section{研究背景与意义}

随着人工智能、传感技术以及高性能计算平台的迅猛发展，自动驾驶技术已经逐渐从概念探索走向实际部署，成为学术界和工业界共同关注的研究
热点。在全球范围内，科技企业、汽车制造商以及研究机构纷纷加大对自动驾驶关键技术的研发投入，旨在构建安全、高效并具备可持续性的智能
交通系统。自动驾驶不仅能够显著降低由人为操作失误所导致的交通事故风险，还能够有效提升道路通行效率、减少交通拥堵，并促进运输能源
结构的优化，从而为缓解城市交通压力、节能减排和环境保护提供重要的技术手段。同时，伴随智慧城市建设的推进，自动驾驶车辆在共享出行
、智慧物流、无人物流仓储等领域的应用也展现出广阔前景。

近年来，自动驾驶系统架构经历了从传统模块化方法向深度学习驱动的端到端方法的发展趋势。端到端控制方案通过深度神经网络直接将多模态
传感器输入（如摄像头、激光雷达、毫米波雷达、惯性传感器等）映射到转向角、加速度、刹车等控制命令，从而避免了传统方案中感知、定
位、规划和控制等多个模块串联可能带来的误差累积问题\cite{end_to_end}。在一定程度上，端到端方法具有更强的特征提取和复杂场景处
理能力，特别是在非结构化道路环境或大规模数据可用性较强的应用场景中表现出显著优势。然而，端到端方法作为典型的数据驱动黑盒模型，
其决策过程缺乏清晰的物理约束与可解释性，使其在安全性、可靠性和可验证性方面难以满足自动驾驶实际部署中的严格要求。尤其是在面对
极端场景、传感器异常或环境分布变化等情况时，端到端控制器可能出现不可控的决策偏差。此外，完全依赖数据驱动的方式往往无法保证系
统稳定性，不利于满足车辆控制领域对性能和安全约束的严格要求。因此，目前端到端方法尚难以完全取代具备可解释性、可分析性和安全保
证的传统控制架构。

在传统自动驾驶控制体系中，轨迹跟踪控制是最核心的关键技术之一，其目标是通过调节车辆控制输入，使车辆能够沿给定参考轨迹行驶，并最小化
横向误差、纵向误差及航向角误差等关键指标\cite{model_and_control_survey}。轨迹跟踪精度直接决定车辆行驶的稳定性、舒适性与安全
性，是连接上层规划模块与底层动力学执行模块的关键环节。然而，轨迹跟踪控制的实现难点在于：车辆动力学模型具有明显的非线性特征，其
特性会随车辆速度、路面状况、载荷变化等多种因素发生显著变化；控制链路中存在的各类延迟、可能的硬件损伤（如车轮执行力下降等驱
动器异常）以及路面摩擦系数变化等难以预先建模的外部扰动，均会导致名义模型与实际系统之间产生不可忽视的偏差。这使得传统控制
器（如PID控制、LQR控制等）在应对复杂、不确定及动态多变工况时难以保持理想性能。尽管模型预测控制（Model Predictive Control, 
MPC）等先进控制方法能够显式处理系统约束并具备一定鲁棒性，但其性能仍高度依赖系统模型的准确性，实际应用中的模型失配问题可能导致
控制性能显著下降。

在此背景下，融合学习方法与控制理论的智能控制技术逐渐成为自动驾驶控制领域的重要发展方向。智能控制通过引入机器学习、概率建模或强
化学习等方法，使控制器具备环境认知、模型补偿甚至策略自学习的能力，从而提升系统在未知、非线性和不确定环境中的适应性。其中，高斯
过程回归（Gaussian Process Regression, GPR）以其严格的概率框架天然支持不确定性量化，能够在学习系统残差动力学的同时提供置信
区间估计，为基于模型的控制器提供有效的在线补偿。相关研究表明，在保证一定可解释性和稳定性的前提下引入数据驱动的学习能力，有望弥补传统方法在
非线性建模方面的不足，同时有效缓解端到端方案在安全性与可信性方面的缺陷。然而，将标准高斯过程回归应用于实时控制场景仍面临计算瓶颈：
其训练与推断复杂度随数据规模呈三次方增长，难以满足在线更新的实时性要求。如何在稀疏近似框架下实现高效、持续的在线模型学习，并将
其无缝集成至模型预测控制中，仍是一个具有挑战性的开放问题。

综上所述，如何设计一种兼具自适应性、鲁棒性与高精度的轨迹跟踪控制方法，以有效应对真实部署中的模型失配、执行器异常与环境扰动，是
自动驾驶领域长期面临的重要研究挑战。

\section{国内外研究现状}

随着控制理论与机器学习技术的快速发展，模型预测控制在复杂、非线性以及存在不确定性的系统中得到了广泛关注。然而，传统MPC依赖精确
的系统模型，难以应对实际工程中的建模误差、扰动以及环境变化。因此，将学习方法引入MPC以提升预测性能与鲁棒性成为近年来的研究热点。
本节将从基于学习的模型预测控制、基于高斯过程回归的模型预测控制以及可扩展高斯过程回归及其在线学习三个方面综述当前国内外研究进展。

\subsection{基于学习的模型预测控制}

传统模型预测控制的核心假设是系统动力学模型已知且能够准确描述真实系统。然而在实际工程应用中，由于模型简化、强非线性效应、未建模
动态以及时变环境等因素，获得精确模型往往不现实。传统建模方法主要包括基于物理机理的建模和基于系统辨识的参数化建模，但这两类方法
在复杂场景下均存在局限：物理建模可能无法覆盖系统的全部动态行为，而系统辨识方法的效果则高度依赖激励信号和数据质量。

随着数据驱动技术的发展，基于学习的MPC（Learning-based MPC, LB-MPC）逐渐成为研究热点。其核心思想是在传统名义模型基础上引入
残差模型或补偿模型，利用离线或在线数据修正预测误差，从而提升系统的建模精度和预测能力\cite{wu2019ml_mpc_part2}。

LB-MPC可以分为基于鲁棒MPC（Robust MPC， RMPC）和基于随机MPC（Stochastic MPC, SMPC）的两类，前者确保闭环系统在所有可能的不确定
性场景下仍能保证稳定性和约束满足性
\cite{tanaskovic2014adaptive_rhc, calliess2014conservative_decision, beliakov2006lipschitz_interpolation}
；后者使用概率分布来表征不确定性，约束通常以概率形式表达，其优点是更符合实际任务中噪声和扰动的统计特性，预测结果更加平滑、更贴近物
理特性，同时也存在需对状态分布进行传播（而非仅传播均值），预测步长越长概率分布之间的耦合越复杂，计算复杂度很高等难题和挑战
\cite{soloperto2018lb_rmpc_state_dep_uncertainty}。
\autoref{fig:rmpc_smpc} 展示了 RMPC 与 SMPC 的建模方式对比，可以看出 SMPC 尤其是基于 GP 的方法能够提供更加平滑且带不确定性估
计的动力学模型，更适合复杂动态场景下的预测控制任务。

学习模型的选择多样：可以使用深度神经网络根据新数据在线更新模型参数\cite{moumouh2020robust_mpc_ann}；通过强化学习实现与环境
的实时交互建模\cite{shan2020rl_path_tracking}；
或者使用高斯过程对残差动力学的均值与方差进行建模，从而得到更符合物理特性的概率预测模型。特别是高斯过程模型所提供的不确定性量化能
力，使得LB-MPC在安全性和鲁棒性方面具有显著优势\cite{he_2020_cautious_mpc}。


除了预测模型的学习，LB-MPC还可以在控制器参数\cite{sferrazza2020lb_parameterized_mpc,
brochu2010bayesopt_portfolio, piga2019performance_model_learning_mpc}、终端约束及代价函数
\cite{blanchini2005relatively_optimal_control, brunner2014terminal_set_enlargement, rosolia2018learning_mpc_iterative, 
liu2014safe_set_control, rosolia2019autonomous_racing_lmpc}等模块中引入学习机制。
通过贝叶斯优化可实现MPC权重矩阵的自动调参\cite{marco2016gp_lqr_tuning}，逆优化控制方法则通过人类演示数据学习最优代价函数参数
\cite{rokonuzzaman2020lb_mpc_path_tracking, abbeel2004apprenticeship_learning, menner2019inverse_learning_motion}。
基于采样安全集的迭代学习方法能够保证轨迹跟踪总成本不随迭代增加，并在自动赛
车等高动态任务中取得了成功应用\cite{rosolia2019autonomous_racing_lmpc}。此外，分层强化学习\cite{zuo2012hier_rl_wmr_tracking}
和近端策略优化\cite{chen2020drl_path_tracking, wikipedia_ppo, chen2019adaptive_rl_path_tracking}也被用于移动机器人路径跟踪
，在跟踪误差和控制平滑度之间取得平衡。

\begin{figure}
    \centering
    \includegraphics[width=.75\linewidth]{survey/rmpc_smpc.png}
    \caption{\label{fig:rmpc_smpc}鲁棒控制（左）和随机控制（右）的建模对比\cite{survey_lb_mpc}}
\end{figure}

为了在学习过程中保证安全性，模型预测安全过滤器（Model Predictive Safety Filter）被引入LB-MPC框架
\cite{wabersich2021predictive_safety_filter}。其核心思想是将学习控制器的输出视为期望控制信号，通过
求解一个附加的MPC优化问题对其进行约束修正：若期望控制满足安全约束则直接执行，否则输出满足约束的最近邻控制
信号。该机制的优点在于通用性强，任何基于学习的控制器均可与安全过滤器组合使用，从而在不牺牲学习性能的前提下提升
系统可靠性\cite{tearle2021safety_filter_racing, berkenkamp2017safe_mbrl, zhang2022survey_lb_mpc}。

LB-MPC方法的主要优势在于：能够在复杂非线性系统中获得比名义模型更高的预测精度，同时可对模型误差和外界扰动进行
显式补偿，并通过不确定性建模与安全过滤器增强鲁棒性。其主要挑战则包括：学习模型引入的不确定性可能影响闭环稳定性，
在线计算成本随数据量增加，以及数据质量对预测精度的制约。

\subsection{基于高斯过程回归的模型预测控制}

高斯过程回归与模型预测控制的结合为非线性系统提供了概率预测与不确定性量化能力，近年来受到广泛关注。在GP-MPC框架中，系统动力
学被建模为名义模型与GP学习残差之和，GP模型在提供预测均值的同时输出预测方差，从而为控制器提供对模型误差的概率描述。

不确定性传播是GP-MPC的核心技术挑战之一。由于GP的输出分布依赖于输入，而MPC的多步预测中每一步的状态输入又来自上一步的预测输出，不
确定性会沿预测时域逐步累积和传播。处理这一问题的主要策略分为直接法和间接法两类：直接法为每个预测步长独立训练GP模型，预测分布精确
但训练成本随步长线性增加\cite{Hachino_2007, Hachino_2011, Pfefferkorn_2022}
；间接法则采用递归的一步预测模型进行迭代传播，训练简单但预测分布随迭代变为非高斯\cite{Farmer1989, Hewing2020}。
为近似多步预测的均值与
方差，常用方法包括一阶线性化\cite{Girard_2002}、矩匹配\cite{Opper_1999}、Sigma点传播\cite{Hardy_2014, Ostafew_2016}
以及蒙特卡洛积分等。
也有研究通过引入Lipschitz上界构建多步预测的置信超椭球，用于安全关键系统中的约束紧缩\cite{Koller2018}。

\begin{figure}
    \centering
    \includegraphics[width=.75\linewidth]{survey/gpmpc_overall.png}
    \caption{\label{fig:overall_show}GP-MPC 的建模演示，绿色圆圈为每个预测点的$1-\sigma$不确定度}
\end{figure}

闭环安全性保证是GP-MPC走向实用的关键。由于GP预测带有概率性质，系统约束通常以概率约束（chance constraint）的形式表达，
要求约束在一定概率
水平上得到满足。现有方法大致分为三类。第一类基于有界支持假设，通过状态依赖的约束收紧实现鲁棒可行性
\cite{Bonzanini2022, soloperto2018lb_rmpc_state_dep_uncertainty}。第二
类基于概率意义鲁棒设计，构建高置信度模型误差集合，并在此集合上设计鲁棒控制器
\cite{Srinivas2012, Chowdhury2017, berkenkamp2017safe_mbrl}。第三类为基于采样
的方法，通过对GP进行多条轨迹采样来估计多步联合分布，再基于样本进行约束紧缩，可减少保守性但计算开销较大
\cite{Bradford2019, Bradford2020, Ma2023}。
部分研究进一步引入Lyapunov函数或均方稳定性分析，将概率性质直接整合进学习阶段，为闭环系统提供更严格的理论保证
\cite{Gruner2022, Matschek2023, Faulwasser2016, Beckers2016, Umlauft2017}。
此外，可在训练阶段引入输入不确定性建模，例如变分高斯过程或高斯混合输入，进一步提高预测精度 \cite{McHutchon2011, Damianou2015}。

在实际应用方面，GP-MPC已在自动驾驶、机器人操作和四旋翼控制等任务中取得了积极成果。Ostafew等人\cite{ostafew2016robust_lbm_nmpc}
将GP-MPC应用于户外移动机器人的鲁棒轨迹跟踪，采用Sigma点方法传播预测不确定性；He等人\cite{he_2020_cautious_mpc}在自动赛车任务中
构建GP-MPC框架，利用残差模型估计不确定性并计算安全裕度；Lahr等人\cite{lahr_2023_zero_order}则使用序列二次规划求解GP-MPC问题，
并证明其可收敛至原非线性问题的最优解，为实时求解提供了新的理论依据。

尽管GP-MPC在上述场景中展示了良好的性能，但现有研究大多假设GP模型在控制器运行前已完成离线训练，难以应对系统特性随时间变化或部署后
出现未知扰动的情形。如何在控制执行过程中持续更新GP模型，是GP-MPC走向实用部署的重要开放问题。

\subsection{可扩展高斯过程回归：稀疏方法与在线学习}

标准GP推断的计算复杂度随训练样本数呈三次方增长，存储开销随样本数平方增长，这使得GP在大规模数据和在线学习场景中的直接应用面临严
重的计算瓶颈。为此，研究者从稀疏近似和在线/流式学习两个方向提出了大量可扩展方法。

\textbf{稀疏近似}方法的核心思想是以少量代表性信息替代全部训练数据进行推断，从而将训练和预测的计算复杂度从$\mathcal{O}(N^3)$
降至$\mathcal{O}(NM^2)$，其中$M \ll N$为代表点数量。数据子集方法直接从训练集中选取$M$个样本，可随机选取或基于信息论准则
选取\cite{gonzalez1985clustering, feder1988optimal}，实现简单但可能丢失重要信息。专
家模型方法将数据划分为多个子集并分别训练局部GP，通过Mixture-of-Experts (MoE)\cite{jacobs1991adaptive, tresp2000mixtures, yuksel2012twenty}
或Product-of-Experts (PoE) \cite{hinton2002training, tresp2000bayesian}进行预测融
合，适合并行计算和非平稳数据建模。有限维核近似方法则通过随机傅里叶特征\cite{lazaro2010sparse}
或协方差级数展开将核函数映射到有限维空间，实现线性复杂度的近似推断，例如稀疏谱高斯过程（Sparse Spectrum GP, SSGP）。

\textbf{诱导点}方法是目前最主流的稀疏GP框架。其核心思想是引入$M
$个诱导点作为充分统计量，通过条件独立性假设将训练点与测试点的联合分布分解为对诱导点的条件分布，从而大幅降低计算成本。
代表性方法包括：SoR（Subset of Regressors）\cite{silverman1985spline} 和DTC（Deterministic Training Conditional）
\cite{csato2002sparse}，
对核矩阵进行显式低秩近似；FITC（Fully Independent Training Conditional）\cite{snelson2007local}，在SoR基础上保留
对角修正项以改善协方差结构；以及VFE（Variational Free Energy）\cite{titsias2008variational}，通过最小化变分后验与真
实后验之间的KL散度得到更原则性的近似框架。其中VFE方法具有清晰的变分下界解释，可通过梯度优化同时学习诱导点位置和超参数，已
成为稀疏GP的主流基础框架。不同高斯过程回归方法的计算效率见\autoref{tab:compute_complex}。

高效实现上述稀疏方法可借助以下数值技术：Nyström 近似利用选定列对 Gram 矩阵进行低秩近似，复杂度 $\mathcal{O}(NM^2)$;
共轭梯度（Conjugate Gradient, CG）用于求解线性系统及对数行列式的迭代方法，复杂度为 $\mathcal{O}(PN^2)$，其中 $P \ll N$;
并行化,基于块独立性与矩阵-向量乘（MVM）求解器，可实现 GPU 加速。

\begin{table}[ht]
\centering
\caption{高斯过程回归方法的计算复杂度}
\label{tab:compute_complex}
\begin{tabular}{lcccccc}
\toprule
& GP & SoD & Expert & DTC/FITC & SSGP & VFE \\
\midrule
训练  & $\mathcal O(N^3)$ & $\mathcal O(M^3)$ & $\mathcal O(NM_e^2)$ & $\mathcal O(NM^2)$ & $\mathcal O(Np^2)$ & $\mathcal O(NM^2)$ \\
推理 & $\mathcal O(N^2)$ & $\mathcal O(M^2)$ & $\mathcal O(M_e^2)$  & $\mathcal O(M^2)$  & $\mathcal O(p^2)$  & $\mathcal O(M^2)$ \\
\bottomrule
\end{tabular}
\end{table}

\textbf{在线与流式学习}方法致力于在数据按序列到达的场景中实现GP的增量更新，是将GP集成到实时控制系统的关键技术。

活跃数据选择与字典维护方法通过维护一个固定大小的代表性数据集来控制模型规模。其基本思路是对每个新到达的数据点评估其信
息增益，若超过阈值则纳入字典并移除贡献最小的旧点，否则直接丢弃。常用的信息度量包括基于RKHS的投影误差\cite{petelin_2011_cs}、
KL散度以及后验方差减少量等。该类方法实现简单、计算开销可控，但字典固定大小的约束限制了模型对新动态的适应能力，且数据丢弃策略可
能导致历史信息的不可逆损失。

\textbf{在线诱导点}方法将稀疏GP的诱导点框架扩展到在线设置，通过迭代方式更新诱导点位置及其关联的后验统计量。
在诱导点自适应管理方面，部分工作探索
了根据预测不确定性或误差动态增减诱导点数量的策略，以在模型容量和计算开销之间取得平衡。

Bijl等人\cite{bijl_2015_online}提出了基于在线FITC和PITC近似的序贯更新算法，该工作在FITC和PITC两种稀疏近似框架下，
推导了后验统计量的闭式递推更新公式，在每次观测后通过秩一更新对诱导点变分分布进行修正，
使得模型在新数据到达时能够以较低的计算代价完成更新，无需求解完整的变分优化问题。
该方法在计算效率上具有明显优势，尤其适合高频数据场景，但其近似质量相对于变分方法缺乏理论上界保证。

Gómez-Bombarelli等人提出的Adaptive Sparse Gaussian Process（ASGP）\cite{gomez_asgp}
通过在线评估每个诱导点对模型预测质量的边际贡献，动态决定是否增加新的诱导点或删除冗余点，从而在
模型精度和计算效率之间实现自动权衡。同时通过加入遗忘参数，削减旧数据对系统的影响，以实时适应新数据分布。
与固定字典方法相比，ASGP能够根据数据分布的变化自适应调整模型结构，在非平稳和分布偏移
场景中表现出更好的适应性，但其遗忘特性导致对系统中不变特性的预测精度下降，缺乏合理的记忆特性。

流式变分推断方法从贝叶斯学习的角度为在线GP更新提供了更具原则性的理论框架。Hensman\cite{Hensman_2015}以及Bui等人
\cite{Bui_2017_ssgp}提出了基于变分推断
的流式GP学习算法，将在线学习问题转化为对变分后验参数的序贯更新问题。具体而言，该方法将每次新数据的到来视为对变分下界的局部
优化，通过自然梯度更新诱导点处的变分均值和协方差，实现了在不存储历史数据的前提下对后验分布的持续修正。与基于矩匹配的在线方法
相比，变分框架在理论上提供了更严格的近似质量保证，并天然支持小批量（mini-batch）更新，计算效率相对完整更新方案更高，
但计算复杂度仍无法降低到在实时系统、边缘设备上部署。

在控制系统的在线GP应用方面，部分研究已尝试将在线GP学习与MPC直接结合。Ostafew等人将在线GP残差学习嵌入非线性MPC，实现了移动
机器人在户外地形中的自适应轨迹跟踪\cite{ostafew2016robust_lbm_nmpc}；局部在线GP方法\cite{park_2013_lbrc}则通过维护多
个局部模型并动态分配新数据，在保证局部精度的同时控制了全局计算成本;
Meier等人\cite{Meier_2016_drift, Meier_2015_drift}针对非平稳动力学场景提出了自适应在线GP方法，通
过引入遗忘机制使模型能够追踪系统特性的漂移变化，并在机器人力控制任务中验证了其有效性。

然而，上述工作要么在在线更新机制与变分近似框架的结合深
度上有所欠缺，诱导点的自适应管理与不确定性在控制决策中的有效利用仍有较大改进空间，要么计算复杂度较高，难以在边缘设备中
高效运行。

综上所述，将高斯过程回归与模型预测控制相结合的研究已取得显著进展，但现有工作在在线学习效率、变分近似框架下的诱导点自适应更新，
以及将模型不确定性融入控制代价等方面仍存在不足。特别是针对差速驱动移动机器人在控制延迟、执行器损伤等典型模型失配场景下的在线
自适应轨迹跟踪问题，尚缺乏系统性的研究。本文在上述工作的基础上，提出一种基于变分稀疏高斯过程的在线学习模型预测控制方法，旨在
填补这一空白。

\section{论文内容与章节安排}
本文针对差速驱动移动机器人轨迹跟踪控制中的模型失配问题，提出一种将在线稀疏高斯过程回归与模型预测控制相结合的自适应控制框架。
该框架能够离线和在线学习名义模型的残差动力学，补偿由控制延迟、执行器异常及环境扰动等因素引起的模型偏差，从而
提升轨迹跟踪精度与系统鲁棒性。全文共五章，各章内容安排如下：

第一章为绪论。首先阐述自动驾驶与智能控制的研究背景与现实意义，分析端到端控制方法与传统基于模型控制方法各自的优势与局限，
指出轨迹跟踪控制在复杂工况下面临的模型失配挑战；随后对基于学习的模型预测控制、基于高斯过程回归的模型预测控制
以及可扩展高斯过程回归三个方向的
国内外研究现状进行综述，梳理现有方法的不足与开放问题；最后说明本文的研究内容，介绍各章内容安排。

第二章介绍本文所涉及的理论背景与基础知识。首先建立差速驱动移动机器人的运动学模型，分析其状态空间表示与控制输入结构，
介绍模型预测控制的基本框架，包括滚动优化机制、代价函数设计以及约束处理方法；
随后介绍标准高斯过程回归的基本原理，包括先验建模、核函数选取、后验推断与超参数优化等核心内容，为后续章节的算法设计奠定理论基础。

第三章提出面向在线学习的稀疏高斯过程回归算法。首先介绍稀疏高斯过程的DTC、FITC近似结构，并引入变分推断框架对诱导点后验分布进
行建模，以诱导点处的变分均值与协方差作为模型的充分统计量；在此基础上，设计支持在线更新的变分参数递推机制，使模型能够
随新观测数据持续修正后验估计；进一步提出诱导点的增量扩充策略，以在不重新训练的前提下提升模型对新区域的表达能力。本章
算法以学习系统名义模型的残差动力学为目标，在提供点预测的同时输出预测不确定性，为第四章的控制集成提供接口。

第四章将线稀疏高斯过程回归算法与模型预测控制相结合，构建完整的学习型轨迹跟踪GP-MPC框架。首先介绍高斯过程预测
值在MPC预测时域内的均值与方差传播方法，讨论如何将残差补偿纳入状态预测模型；随后阐述控制框架的整体结构，包括离线初始化
与在线自适应两种工作模式的切换机制，以及不确定性信息在代价函数中的利用方式。最后通过仿真与实物实验对所提方法进行验证：
设计控制延迟与单侧车轮执行力损失两类典型模型失配场景，与基准MPC控制器进行对比，从轨迹跟踪精度与模型在线补偿效果两个维度
评估所提框架的有效性。

第五章为总结与展望。对全文的主要研究工作与结论进行归纳总结，分析所提方法在计算效率、模型表达能力及适用场景等方面存在
的局限性，并就诱导点自适应管理、多步不确定性传播、以及向更复杂平台的迁移应用等方向提出未来研究展望。

各章节之间的逻辑关系如图1.X所示。第二章为全文提供理论基础；第三章在此基础上提出在线学习的稀疏高斯过程回归算法；
第四章设计了可将GP算法与MPC框架深度集成的方案并通过实验加以验证；第一章与第五章分别从宏观层面完成研究动机的阐释与研究结论的收束，
共同构成完整的研究闭环。